% Sección 06 - Protocolo experimental (Dayver)
% Cada bloque "% DATOS:" es lo que tienes que contar con tus palabras.
% Borra los comentarios cuando escribas el texto.

\subsection{Entorno}
% DATOS: la corrida oficial se hizo el 22/11/2026 a las 11:40 en la PC de Mario
% (equipo ReyShadow). Los nodos corrieron como procesos en localhost, puertos
% 5000 a 5003, sin Docker. Sockets y gRPC nunca corrieron a la vez.
% Las condiciones las registró src/experimento/condiciones.py y están en la
% tabla de abajo (no las copies al texto, solo menciona lo principal:
% Windows 11, Intel Core i7-14700, 63,7 GB de RAM, Python 3.13.15).

\begin{table}[H]
  \centering
  \caption{Condiciones de la máquina de la corrida oficial}
  \label{tab:condiciones}
  % Generado por analysis/analizar.py. No editar a mano.
% condiciones_mcastro.csv
\begin{tabular}{l l}
\toprule
Campo & Valor \\
\midrule
fecha\_hora & 2026-11-22T11:40:39 \\
equipo & ReyShadow \\
sistema\_operativo & Windows 11 \\
version\_so & 10.0.26300 \\
arquitectura & AMD64 \\
procesador & Intel(R) Core(TM) i7-14700 \\
nucleos\_logicos & 28 \\
ram\_gb & 63.7 \\
python & 3.13.15 \\
implementacion\_python & CPython \\
modo\_nodos & procesos en localhost (puertos 5000, 5001, 5002) \\
version\_grpcio & 1.85.3 \\
version\_grpcio-tools & 1.85.3 \\
version\_protobuf & 7.36.2 \\
version\_pandas & 3.0.2 \\
version\_matplotlib & 3.11.2 \\
version\_jinja2 & 3.1.6 \\
\bottomrule
\end{tabular}

\end{table}

\subsection{Corridas}
% DATOS: un solo comando corre todo:
%   python -m src.experimento.correr_experimento
% Son 78 corridas:
%  - Latencia, 48 corridas: 2 escenarios (secuencial: nodo-B -> nodo-A;
%    concurrente: nodo-B, nodo-C y nodo-D a la vez -> nodo-A)
%    x 2 tamaños (100 B y 10 240 B) x 6 repeticiones (r0 a r5) x 2 mecanismos.
%    Cada cliente hace 100 envíos por corrida y mide el RTT con time.perf_counter().
%  - Lamport, 30 corridas: 2, 3 y 4 nodos x 5 repeticiones x 2 mecanismos,
%    20 mensajes en total por corrida en la ronda de dos fichas.
% Cada corrida tiene una etiqueta (s100-r1, c10240-r3, n3-r2...) que reciben
% todos sus nodos con --corrida; así los CSV de una corrida se emparejan.

\subsection{Calentamiento y repeticiones}
% DATOS: r0 es la corrida de calentamiento: se ejecuta igual que las demás
% pero el análisis la descarta. Por qué: la primera vez se cargan módulos de
% Python y gRPC, se crean hilos y se llenan cachés, y eso infla los tiempos.
% Después vienen 5 repeticiones medidas (r1 a r5). La unidad de análisis es
% cada repetición (ver el criterio de decisión de la sección 2).
% Lamport no tiene calentamiento porque ahí se verifica corrección, no tiempo.

\subsection{Orden de los mecanismos}
% DATOS: en cada repetición los dos mecanismos corren uno detrás del otro y se
% alterna cuál va primero: r0 sockets-gRPC, r1 gRPC-sockets, r2 sockets-gRPC...
% Por qué: si siempre fuera primero el mismo, podría sacar ventaja o desventaja
% por el estado del sistema (cachés, otros procesos, temperatura del CPU).
% No queda perfectamente balanceado: con 5 repeticiones, gRPC va primero en 3
% (r1, r3, r5) y sockets en 2 (r2, r4).

\subsection{Control de las corridas}
% DATOS: antes de cada corrida el script revisa que los puertos estén libres y
% arranca primero el servidor. Al terminar revisa que cada nodo haya dejado su
% CSV de eventos y cada cliente su CSV con 100 envíos; si algo falta, se detiene.
% Deja un resumen de las 78 corridas en data/corridas_mcastro_20261122-114039.csv,
% junta la latencia en data/latency_sockets.csv y data/latency_grpc.csv, y
% los registros de pantalla en logs/. Después de la corrida se generó
% data/SHA256SUMS para poder comprobar que los CSV no se modificaron.
